% Lösungen Einheit 2
\einheitenkopf[le2][2 p-q-Formel]{Einheit 2 von 4 · Normalform und p-q-Formel}
\erg{24}{a) p-q-Formel \quad b) Wurzelziehen \quad c) rückwärts rechnen \quad d) Wurzelziehen \quad e) p-q-Formel (erst ordnen)}
\erg{25}{a) p = 7, q = 2 \quad b) p = −2, q = 9 \quad c) p = 5, q = −1 \quad d) p = −1, q = −20 \quad e) p = 0, q = −11}
\erg{26}{a) nicht teilen \quad b) teilen durch 2 \quad c) teilen durch −1 \quad d) teilen durch 0,5 \quad e) nicht teilen}
\erg{27}{a) −3 $\pm$ $\surd$(9 − 8) = −3 $\pm$ 1; $x_1$ = −2, $x_2$ = −4 \quad b) −5 $\pm$ 3; $x_1$ = −2, $x_2$ = −8 \quad c) −3 $\pm$ 2; $x_1$ = −1, $x_2$ = −5 \quad d) −5 $\pm$ 2; $x_1$ = −3, $x_2$ = −7}
\erg{28}{a) 2 $\pm$ 4; $x_1$ = 6, $x_2$ = −2 \quad b) −1,5 $\pm$ $\surd$12,25 = −1,5 $\pm$ 3,5; $x_1$ = 2, $x_2$ = −5 \quad c) 2 $\pm$ $\surd$3; $x_1$ $\approx$ 3,73, $x_2$ $\approx$ 0,27 \quad d) −0,5 $\pm$ $\surd$3,25; $x_1$ $\approx$ 1,30, $x_2$ $\approx$ −2,30 \quad e) x² − 8x + 14 = 0; 4 $\pm$ $\surd$2; $x_1$ $\approx$ 5,41, $x_2$ $\approx$ 2,59}
\erg{29}{a) x² + 2x − 8 = 0; $x_1$ = 2, $x_2$ = −4 \quad b) x² − 6x + 5 = 0; $x_1$ = 5, $x_2$ = 1 \quad c) x² − 5x + 4 = 0; $x_1$ = 4, $x_2$ = 1 \quad d) x² + 6x − 16 = 0; $x_1$ = 2, $x_2$ = −8}
\erg{30}{a) x² − 2x − 3 = 0; $x_1$ = 3, $x_2$ = −1 \quad b) x² + 3x − 4 = 0; $x_1$ = 1, $x_2$ = −4 \quad c) x² − 6x + 8 = 0; $x_1$ = 4, $x_2$ = 2 \quad d) x² − 2x − 8 = 0; $x_1$ = 4, $x_2$ = −2 \quad e) 3x² − 3x − 36 = 0, geteilt durch 3: x² − x − 12 = 0; $x_1$ = 4, $x_2$ = −3}
\erg{31}{a) x² + 4x − 12 = 0; $x_1$ = 2, $x_2$ = −6 \quad b) x² + 2x + 1 = 3x + 7, x² − x − 6 = 0; $x_1$ = 3, $x_2$ = −2 \quad c) x² − 6x + 8 = −2x + 8, x² − 4x = 0; $x_1$ = 4, $x_2$ = 0 \quad d) −x² + 6x − 1 = −8, −x² + 6x + 7 = 0, x² − 6x − 7 = 0; $x_1$ = 7, $x_2$ = −1}
\erg{32}{a) x = −3 \quad b) x = 5 \quad c) x = 0,5 \quad d) x² + 3x + 2,25 = 0; x = −1,5}
\erg{33}{a) 1 − 5 = −4: keine Lösung \quad b) 9 − 10 = −1: keine Lösung \quad c) 0,25 − 1 = −0,75: keine Lösung \quad d) x² + 2x + 3,5 = 0; 1 − 3,5 = −2,5: keine Lösung}
\erg{34}{a) 3, zwei \quad b) 0, eine \quad c) −3, keine \quad d) 1,25, zwei \quad e) −0,75, keine \quad f) 2, zwei}
\erg{35}{a) 4 + 4 − 8 = 0, wA \quad b) 1 + 5 + 4 = 10 $\neq$ 0, fA \quad c) 1 + 2 − 3 = 0, wA \quad d) 9 + 3 − 6 = 6 $\neq$ 0, fA \quad e) 16 − 12 − 4 = 0, wA}
\erg{36}{a) Wurzelziehen: x² = 25; $x_1$ = 5, $x_2$ = −5 \quad b) p-q-Formel: $x_1$ = 4, $x_2$ = −6 \quad c) rückwärts: x + 3 = 4 oder x + 3 = −4; $x_1$ = 1, $x_2$ = −7 \quad d) ordnen, p-q-Formel: x² − 6x − 16 = 0; $x_1$ = 8, $x_2$ = −2}
\erg{37}{a) Nicht durch 2 geteilt; x² + 8x + 7 = 0, −4 $\pm$ 3; $x_1$ = −1, $x_2$ = −7 \quad b) −p/2 ist 5, nicht −5; 5 $\pm$ 4; $x_1$ = 9, $x_2$ = 1 \quad c) Unter der Wurzel steht minus q: 16 − 12 = 4; 4 $\pm$ 2; $x_1$ = 6, $x_2$ = 2}
\erg{38}{a) Vor x² steht 2; die Formel gilt nur für die Normalform mit 1 vor x², also erst durch 2 teilen. \quad b) Aus einer negativen Zahl gibt es keine Wurzel, also gibt es keine Zahl x, die passt.}
\erg{39}{a) x² − 2x − 1 = x + 3, x² − 3x − 4 = 0; $x_1$ = 4, $x_2$ = −1 \quad b) y = 7 und y = 2; S$_1$(4 | 7), S$_2$(−1 | 2) \quad c) Die Graphen schneiden sich in (4 | 7) und (−1 | 2).}
